\newpage
\section{Duração e Cronograma}

\subsection{Análise da Duração do Projeto}

O desenvolvimento do sistema \textbf{Tropical Mix} foi estruturado em sete fases principais, abrangendo desde o planejamento e definição técnica até a entrega da aplicação final. O cronograma foi planejado para garantir um avanço progressivo, permitindo a validação de cada etapa antes do início da seguinte.  

\subsubsection{Fase 1 – Planejamento e Desenho da Aplicação (23/09/2025 – 24/10/2025)}

Esta fase abrangeu o levantamento de requisitos, definição das tecnologias e elaboração da arquitetura da aplicação. Também foram desenvolvidos os diagramas de modelagem, o Modelo Entidade-Relacionamento (MER) e o Diagrama Entidade-Relacionamento (DER), além da documentação inicial sobre regras de negócio e segurança da informação.  

Durante esta etapa, também foram realizadas a configuração do ambiente de desenvolvimento, organização da estrutura do projeto e implementação inicial da Prova de Conceito (POC).  

\subsubsection{Fase 2 – Desenvolvimento da Base do Sistema (25/10/2025 – 21/11/2025)}

Nesta fase, foram desenvolvidas as primeiras funcionalidades estruturais do sistema, incluindo criação das tabelas do banco de dados, implementação das rotas iniciais da API e integração com o banco MySQL.  

Também foi iniciada a documentação técnica do projeto e a estruturação inicial das telas da aplicação.  

\subsubsection{Fase 3 – Produto Mínimo Viável (MVP) (22/11/2025 – 05/12/2025)}

Nesta fase, foi implementada uma versão funcional do sistema, contendo as principais operações: controle de estoque, registro de vendas, autenticação de usuários. 

O MVP serviu como base para avaliações práticas junto ao contexto da loja, possibilitando ajustes antes da entrega final.

Também foram realizadas integrações entre \textit{front-end} e \textit{back-end}, implementação de relatórios e testes iniciais das funcionalidades principais do sistema.

Além disso, foi realizada a hospedagem da aplicação em ambiente AWS, bem como a integração entre API, banco de dados e servidor de hospedagem.

\subsubsection{Fase 4 – Expansão das Funcionalidades Administrativas (06/12/2025 – 31/01/2026)}

Esta etapa foi destinada à implementação de funcionalidades administrativas, incluindo gerenciamento de usuários, permissões de acesso, notificações e cadastro de fornecedores.  

Também foram realizados ajustes estruturais nas tabelas do banco de dados e melhorias na organização das rotas e módulos da aplicação.  

\subsubsection{Fase 5 – Refinamento e Infraestrutura (01/02/2026 – 31/03/2026)}

Nesta fase, foram implementadas melhorias voltadas à usabilidade e infraestrutura do sistema, incluindo filtros, menus \textit{dropdown}, validações de produtos vencidos e ajustes nas funcionalidades de estoque e produtos.    

\subsubsection{Fase 6 – Testes e Validação Final (01/04/2026 – 30/04/2026)}

Esta etapa concentrou-se na realização de testes unitários, testes internos da equipe e testes com usuários, permitindo a identificação e correção de falhas no sistema.  

Também foram realizadas otimizações de desempenho e revisão completa da documentação técnica e acadêmica do projeto.  

\subsubsection{Fase 7 – Entrega Final (01/05/2026 – 27/05/2026)}

A entrega final compreende a versão completa e refinada do sistema, com todas as funcionalidades implementadas, correções aplicadas e testes de desempenho concluídos. Essa etapa também incluiu a revisão dos requisitos de segurança e conformidade com a LGPD, além da documentação técnica e do manual de uso.  

Durante esta fase, foram realizados os ajustes finais do sistema, revisão completa do código-fonte, preparação da apresentação do projeto e entrega oficial da aplicação finalizada.